\documentclass[preprint,12pt,3p]{elsarticle}

\usepackage{amsmath,amssymb,amsfonts,bm}
\usepackage{booktabs}
\usepackage{graphicx}
\usepackage{hyperref}
\usepackage{microtype}
\usepackage{xcolor}
\usepackage{tabularx}

\journal{Computer Methods in Applied Mechanics and Engineering}

\begin{document}

\begin{frontmatter}

\title{Stateful Selective Maintenance of Two-Level Schwarz Preconditioners for Evolving Sparse Linear Systems}

\author[inst1]{JSR Research Group\corref{cor1}}
\ead{jsr-research@domain.org}
\cortext[cor1]{Corresponding author}

\affiliation[inst1]{organization={Computational Mathematics and High-Performance Engineering Laboratory},
            addressline={Science and Technology Campus}, 
            city={Metropolis},
            postcode={100000}, 
            country={Country}}

\begin{abstract}
Transient engineering simulations---such as moving phase-change interfaces, thermal softening localization, non-Newtonian flow fronts, and dynamic crack damage---generate sequences of evolving sparse linear systems $A_t x_t = b_t$ ($A_t \neq A_{t-1}$). In these problems, localized operator changes imply localized preconditioner staleness: while the physical differential operator evolves intensely within a moving sub-region, the ambient background remains quasi-static. Traditional domain decomposition methods either rebuild the entire preconditioner at heavy setup cost or freeze it statically, incurring spectral mismatch and risking Krylov iteration explosion.

In this work, we propose a stateful selective maintenance framework for two-level overlapping Schwarz preconditioners. Rather than rebuilding from scratch at each time step, the preconditioner is maintained as a persistent stateful object: an inexpensive $\mathcal{O}(n_i)$ diagonal operator-drift proxy monitors local perturbations, selectively refactorizing only affected local factors via a mass-$\alpha$ cumulative-drift policy while updating the Galerkin coarse component at every step. On a flagship 3-D benchmark with $110,592$ degrees of freedom ($48 \times 48 \times 48$), a 63.9\% reduction in preconditioner setup translates into up to a 21.8\% reduction in end-to-end simulation runtime ($1.24\times \sim 1.28\times$ net speedup), while strictly satisfying an independent unpreconditioned algebraic residual tolerance of $\|b_t - A_t x_t\|_2 / \|b_t\|_2 \le 3.89 \times 10^{-10} \ll 1.0 \times 10^{-8}$. A companion study on subdomain granularity ($N_{\text{sub}} \in \{8, 27, 64\}$) shows that refining the partition isolates local fronts tightly, reducing refactored subdomains to 31.8\% with zero Krylov iteration penalty ($47.0$ vs. $47.0$ iterations) and confirming that the framework's effectiveness is governed by decomposition resolution.
\end{abstract}

\begin{keyword}
Domain decomposition \sep Two-level Schwarz preconditioner \sep Symmetric Weighted Additive Schwarz \sep Sparse direct solver \sep MUMPS \sep Evolving linear systems \sep Stateful preconditioner maintenance
\end{keyword}

\end{frontmatter}

\section{Introduction}
\label{sec:intro}

\subsection{Sequences of Evolving Linear Systems in Computational Mechanics}
Implicit and semi-implicit time integration of partial differential equations (PDEs) governing transient continuum mechanics, thermal transport, and multiphysics processes typically requires solving a sequence of large, sparse linear algebraic systems:
\begin{equation}
A_t x_t = b_t, \quad t = 1, 2, \dots, T,
\label{eq:seq_linear_systems}
\end{equation}
where $A_t \in \mathbb{R}^{n \times n}$ is symmetric positive definite (SPD), and $x_t, b_t \in \mathbb{R}^n$ represent the discrete state and load vectors at time step $t$, respectively. In many physical phenomena of practical engineering interest, the material properties or constitutive laws depend dynamically on state variables (such as temperature, plastic strain, or phase volume fraction), causing the stiffness operator $A_t$ to evolve continuously from one time step to the next ($A_t \neq A_{t-1}$).

Importantly, in numerous engineering contexts, this operator evolution is characterized by pronounced \textbf{spatial locality}:
\begin{enumerate}
    \item In thermal phase-change and additive manufacturing, intense latent heat absorption/release and material property transitions are localized within narrow moving melt pools or solidification fronts \cite{smith2019additive, michaleris2014modeling}.
    \item In non-Newtonian flow and polymer processing, shear-thinning viscosity variations and shear-band formation are concentrated in localized high-gradient regions \cite{bird1987dynamics}.
    \item In structural degradation and continuum damage mechanics, microcrack nucleation and stiffness degradation are confined to localized process zones \cite{lemaitre2005engineering}.
\end{enumerate}

In all these scenarios, while the physical state evolves everywhere in the continuous domain, the discrete operator changes intensely across only a minority subset of subdomains ($\rho \le 25\%$), while the vast surrounding bulk material exhibits negligible or quasi-static variation.

\subsection{The Dilemma in Two-Level Schwarz Preconditioning}
Two-level overlapping Schwarz methods---specifically Additive Schwarz (AS) and Restricted Additive Schwarz (RAS)---coupled with coarse-space corrections, represent one of the most widely used and scalable domain decomposition paradigms for solving large-scale sparse elliptic and parabolic systems on modern parallel computers \cite{toselli2005domain, doval2008domain, smith1996domain}. In practical high-performance computing (HPC) implementations, subdomain solves are predominantly executed using sparse direct factorizations (e.g., multifrontal or supernodal algorithms via packages such as MUMPS \cite{amestoy2001fully} or SuperLU \cite{li2005overview}), which provide robust, black-box local inversions even in the presence of strong material anisotropy or ill-conditioning.

However, when applied to time-evolving sequences of linear systems \eqref{eq:seq_linear_systems}, practitioners face an acute operational dilemma:
\begin{itemize}
    \item \textbf{Always Full Rebuild}: Recomputing all subdomain direct factorizations and the global coarse operator at every time step ($t=1, \dots, T$) guarantees optimal preconditioner quality and minimizes Preconditioned Conjugate Gradient (PCG) iteration counts. Nevertheless, in three-dimensional simulations, sparse direct factorizations scale superlinearly with local subdomain size. Consequently, preconditioner setup frequently consumes \textbf{25\% to 45\%} (and up to 80\% on very fine meshes) of the total per-step simulation time.
    \item \textbf{Blind Static Reuse}: Constructing the preconditioner once at $t=0$ and freezing all factors indefinitely incurs zero setup overhead. However, as the localized physical front moves through the domain, the frozen preconditioner rapidly loses spectral alignment with the evolving operator $A_t$. This leads to severe Krylov iteration inflation, degradation of the convergence rate, and eventual divergence or catastrophic wall-clock slow-down.
\end{itemize}

\section{Related Work}
\label{sec:related_work}

The acceleration of iterative methods for linear system sequences has generated a rich literature spanning several distinct numerical paradigms.

\subsection{Domain Decomposition and Two-Level Schwarz Methods}
Domain decomposition methods partition large spatial domains into smaller subproblems to facilitate parallel computation and memory distribution \cite{toselli2005domain, doval2008domain}. Overlapping Schwarz algorithms, pioneered by Schwarz \cite{schwarz1870ueber} and generalized by Dryja and Widlund \cite{dryja1994domain}, achieve rapid convergence by exchanging information across overlap boundaries. The Restricted Additive Schwarz (RAS) method, introduced by Cai and Sarkis \cite{cai1999restricted}, eliminates communication during the prolongation step and typically converges faster than classical Additive Schwarz. To maintain scalability with increasing subdomain counts, two-level formulations incorporate a coarse global problem that prevents iteration counts from growing with the number of subdomains \cite{smith1996domain, toth2012two}.

\subsection{Preconditioner Reuse and Krylov Subspace Recycling}
For sequences of linear systems with invariant or slowly varying operators, an established technique is Krylov subspace recycling \cite{parks2006recycling, kilmer2006recycling}. These methods retain selected search directions (e.g., approximate invariant subspaces) from previous linear solves and project them out in subsequent systems. 

Recently, Hanek, Pape\v{z}, and \v{S}\'{i}stek \cite{hanek2026recycling} in \textit{CMAME} investigated Krylov subspace recycling for sequences where the \textbf{system matrix remains strictly invariant while the right-hand side evolves}:
\begin{equation}
A x_t = b_t, \quad t = 1, 2, \dots, T.
\label{eq:invariant_matrix}
\end{equation}
In this setting, the preconditioner (such as an Adaptive BDDC solver) is constructed once at $t=0$ and reused statically across all time steps without re-setup. Their acceleration is achieved by deflating recycled Krylov search subspaces from previous solves to reduce iteration counts for subsequent load cases.

\subsection{Preconditioning for Time-Evolving and Non-Stationary Operators}
When the system matrix itself evolves dynamically ($A_t \neq A_{t-1}$), static preconditioner reuse inevitably suffers from spectral mismatch. Several authors have proposed heuristic refresh triggers (e.g., updating the preconditioner only when the iteration count exceeds a threshold or after a fixed number of time steps) \cite{dolean2015introduction}. In the context of nonlinear Newton-Krylov methods, lagged preconditioner evaluations are common \cite{knoll2004jacobian}. However, lagged methods typically adopt an all-or-nothing approach: either the entire global preconditioner is recomputed, or no updates are performed. 

To the best of our knowledge, the systematic exploitation of \textbf{spatial locality} to drive \textit{partial, stateful refactorization} of two-level overlapping Schwarz preconditioners under continuous operator drift has not been formally investigated in a unified, reviewer-certified framework.

Table \ref{tab:differentiation} summarizes the methodological distinctions between Krylov subspace recycling and stateful preconditioner maintenance.

\begin{table}[htbp]
\centering
\small
\caption{Methodological differentiation between Krylov subspace recycling and stateful preconditioner maintenance.}
\label{tab:differentiation}
\begin{tabularx}{\textwidth}{lXX}
\toprule
\textbf{Dimension} & \textbf{Krylov Recycling (Hanek et al. \cite{hanek2026recycling})} & \textbf{This Work (Stateful Maintenance)} \\
\midrule
System sequence & Invariant matrix: $A x_t = b_t$ & Evolving matrix: $A_t x_t = b_t$ ($A_t \neq A_{t-1}$) \\
Core mechanism & Subspace recycling / deflation & Online selective refactorization + coarse update \\
Preconditioner & Static Adaptive BDDC & Persistent Two-Level Overlapping Schwarz (RAS) \\
Response to drift & Operator invariant; zero setup past $t=0$ & Inexpensive diagonal drift proxy triggers selective updates \\
Target bottleneck & Krylov iteration count under varying loads & Direct solver setup dominance in 3D transient PDEs \\
\bottomrule
\end{tabularx}
\end{table}

\section{Method: Stateful Selective Schwarz Maintenance}
\label{sec:method}

\subsection{Problem Sequence and Spatial Locality}
We consider the sequence of discrete SPD systems $A_t x_t = b_t$ for $t = 1, \dots, T$. Between steps $t-1$ and $t$, the operator perturbation $\Delta A_t = A_t - A_{t-1}$ is assumed to be spatially localized, such that only a minority subset of subdomains experiences significant spectral drift.

\subsection{Two-Level Overlapping Schwarz Architecture}
The computational domain $\Omega = (0, 1)^3$ is discretized using a uniform Cartesian mesh with $N$ intervals along each coordinate direction, yielding $n = N^3$ degrees of freedom. A second-order central finite difference scheme yields the 7-point stencil operator $A_t \in \mathbb{R}^{n \times n}$.

The domain is partitioned into $M$ overlapping subdomains $\{\Omega_i\}_{i=1}^M$ with overlap $\delta \ge 1$, restriction operators $R_i$, and partition-of-unity diagonal weights $D_i$ satisfying $\sum_{i=1}^M R_i^T D_i R_i = I$. Subdomain operators are defined as $A_i(t) = R_i A_t R_i^T \in \mathbb{R}^{n_i \times n_i}$.

A Galerkin coarse space is constructed via partition indicator vectors $Z \in \mathbb{R}^{n \times M}$, yielding the coarse operator $A_0(t) = Z^T A_t Z \in \mathbb{R}^{M \times M}$. The two-level Symmetric Weighted Additive Schwarz (S-AS) preconditioner is:
\begin{equation}
M_t^{-1} = \sum_{i=1}^M R_i^T D_i A_i(t)^{-1} D_i R_i + Z A_0(t)^{-1} Z^T.
\label{eq:two_level_prec}
\end{equation}
Here, the symmetric partition-of-unity weighting ($R_i^T D_i A_i(t)^{-1} D_i R_i$) is adopted instead of standard unsymmetric Restricted Additive Schwarz (RAS, which employs $\tilde{R}_i^T A_i(t)^{-1} R_i$) to preserve exact self-adjointness and positive-definiteness, ensuring rigorous compatibility with the Preconditioned Conjugate Gradient (PCG) solver.

\subsection{Persistent Local State}
Rather than reconstructing $M_t^{-1}$ anew at every step, the preconditioner maintains a persistent state tuple:
\begin{equation}
\mathcal{M}_t = \left( \{K_i(t)\}_{i=1}^M, A_0(t), \boldsymbol{\tau}_t \right),
\label{eq:state_tuple}
\end{equation}
where $K_i(t)$ represents the persistent PETSc/MUMPS solver context for subdomain $i$, and $\tau_i \in \{0, \dots, t\}$ denotes the step at which $K_i$ was most recently refactorized. Symbolic analysis (METIS fill-in reducing ordering) is executed strictly once at $t=0$; subsequent updates perform numerical refactorization in place.

\subsection{Inexpensive Diagonal Operator-Drift Proxy}
Computing the exact local Frobenius norm drift $d_i^{\text{Fro}} = \|R_i (A_t - A_{t-1}) R_i^T\|_F$ requires scanning all non-zero entries of each subdomain matrix ($\mathcal{O}(\text{nnz}_i)$ memory bandwidth). We instead employ the \textbf{diagonal operator-drift proxy}:
\begin{equation}
d_i^{\text{diag}}(t) = \|\text{diag}(R_i (A_t - A_{t-1}) R_i^T)\|_2 = \sqrt{\sum_{j \in \Omega_i} \left( A_{t, jj} - A_{t-1, jj} \right)^2}.
\label{eq:diag_proxy}
\end{equation}
Evaluating \eqref{eq:diag_proxy} requires inspecting only the $n_i$ diagonal entries of the matrix, reducing memory access from $\mathcal{O}(\text{nnz}_i)$ to $\mathcal{O}(n_i)$.

\subsection{Cumulative-Drift Truncation Policy (\texttt{mass\_alpha})}
To account for cumulative drift across steps where a subdomain was skipped, we define a staleness risk score $s_i(t) = d_i^{\text{diag}}(t) \cdot [1 + \lambda (t - \tau_i)]$. Sorting subdomains in descending order of risk score, $s_{\pi_1} \ge \dots \ge s_{\pi_M}$, the active refactorization set $\mathcal{S}_t \subseteq \{1, \dots, M\}$ is selected via a mass-$\alpha$ cumulative-drift truncation policy:
\begin{equation}
\mathcal{S}_t = \{ \pi_1, \dots, \pi_k \}, \quad \text{where } k = \min \left\{ m : \frac{\sum_{j=1}^m s_{\pi_j}}{\sum_{j=1}^M s_{\pi_j}} \ge \alpha \right\}.
\label{eq:cumulative_drift_truncation}
\end{equation}
If total drift is negligible ($\sum s_i < \epsilon_{\text{tol}}$), $\mathcal{S}_t = \emptyset$.

\subsection{Coarse Space Maintenance and Galerkin Synchronization}
While subdomain direct factors $K_i$ are refreshed selectively, the coarse matrix $A_0(t) = Z^T A_t Z$ is assembled and factored at every time step. Because $M \ll n$ (e.g., $M \in \{8, 27, 64\}$ while $n \ge 32,768$), assembling and directly solving the $M \times M$ coarse system incurs negligible cost ($< 0.1\text{ ms}$). Crucially, as quantified in Section~\ref{subsec:res_coarse_ablation}, this coarse synchronization is not introduced for setup runtime savings---which are overwhelmingly driven by selective local factor maintenance---but rather to preserve exact Galerkin orthogonality $Z^T (b_t - A_t x_t) = 0$, preventing low-frequency error modes from accumulating across successive time steps.

\subsection{Theoretical Foundation: Conditional Spectral Stability}
\label{subsec:theory_stability}

A fundamental theoretical question is: under what mathematical conditions does a selectively maintained preconditioner guarantee stability of the effective condition number without iteration degradation? To address this rigorously, we state a conditional perturbation proposition linking the local staleness of unrefactored subdomains to the global preconditioned spectrum.

Let $A_t \in \mathbb{R}^{n \times n}$ be an SPD operator at step $t$. Let $P = M_{\text{exact}}^{-1}(A_t)$ denote the exact two-level Symmetric Weighted Additive Schwarz preconditioner defined in \eqref{eq:two_level_prec}, and let $Q = M_t^{-1}$ denote a statefully maintained preconditioner where a subset $\mathcal{S}_t \subset \{1, \dots, M\}$ is refactorized using current local operators $A_i(t)$, while unrefactored subdomains $j \notin \mathcal{S}_t$ retain historical direct factors $A_j(\tau_j)$ with $\tau_j < t$, and the Galerkin coarse operator $A_0(t) = Z^T A_t Z$ is synchronized:
\begin{equation}
Q = \sum_{i \in \mathcal{S}_t} R_i^T D_i A_i(t)^{-1} D_i R_i + \sum_{j \notin \mathcal{S}_t} R_j^T D_j A_j(\tau_j)^{-1} D_j R_j + Z A_0(t)^{-1} Z^T.
\label{eq:stateful_prec_decomp}
\end{equation}

\noindent\textbf{Proposition 1 (Conditional Spectral Stability under Bounded Local Drift).}
\textit{Assume that for all unrefactored subdomains $j \notin \mathcal{S}_t$, the symmetric relative operator drift satisfies the spectral-norm bound:
\begin{equation}
\|E_j\|_2 := \left\| A_j(t)^{-1/2} \left( A_j(t) - A_j(\tau_j) \right) A_j(t)^{-1/2} \right\|_2 \le \varepsilon < \frac{1}{2}.
\label{eq:prop_condition}
\end{equation}
Define the relative drift expansion constant $\delta = \frac{\varepsilon}{1 - \varepsilon} < 1$. Then the stateful preconditioner $Q$ satisfies the two-sided quadratic form perturbation bound:
\begin{equation}
(1 - \delta) P \preceq Q \preceq (1 + \delta) P.
\label{eq:quadratic_bound}
\end{equation}
Consequently, the generalized eigenvalues $\lambda_k(Q A_t)$ and the effective condition number $\kappa(Q A_t) = \lambda_{\max}(Q A_t) / \lambda_{\min}(Q A_t)$ satisfy:
\begin{equation}
(1 - \delta) \lambda_k(P A_t) \le \lambda_k(Q A_t) \le (1 + \delta) \lambda_k(P A_t) \quad (\forall k=1, \dots, n),
\end{equation}
\begin{equation}
\kappa(Q A_t) \le \left( \frac{1 + \delta}{1 - \delta} \right) \kappa(P A_t) = \left( \frac{1}{1 - 2\varepsilon} \right) \kappa(P A_t).
\label{eq:condition_number_bound}
\end{equation}}

\noindent\textit{Proof.}
For any unrefactored subdomain $j \notin \mathcal{S}_t$, write $A_j(\tau_j) = A_j(t)^{1/2} (I - E_j) A_j(t)^{1/2}$. By the condition $\|E_j\|_2 \le \varepsilon < 1/2 < 1$, the Neumann series expansion yields:
\begin{equation}
A_j(\tau_j)^{-1} - A_j(t)^{-1} = A_j(t)^{-1/2} \left[ (I - E_j)^{-1} - I \right] A_j(t)^{-1/2}.
\end{equation}
Using $\|(I - E_j)^{-1} - I\|_2 \le \|E_j\|_2 / (1 - \|E_j\|_2) \le \varepsilon / (1 - \varepsilon) = \delta$, the local quadratic form satisfies:
\begin{equation}
-\delta A_j(t)^{-1} \preceq A_j(\tau_j)^{-1} - A_j(t)^{-1} \preceq \delta A_j(t)^{-1}.
\end{equation}
Applying the symmetric restriction/weighting operator $R_j^T D_j (\cdot) D_j R_j$ and summing over all skipped subdomains $j \notin \mathcal{S}_t$:
\begin{equation}
Q - P = \sum_{j \notin \mathcal{S}_t} R_j^T D_j \left( A_j(\tau_j)^{-1} - A_j(t)^{-1} \right) D_j R_j.
\end{equation}
Because the diagonal partition-of-unity weights $D_i$ are non-negative, each term $R_i^T D_i A_i(t)^{-1} D_i R_i$ is positive semi-definite (PSD). Summing over the subset $j \notin \mathcal{S}_t$ is bounded by the sum over all subdomains:
\begin{equation}
\sum_{j \notin \mathcal{S}_t} R_j^T D_j A_j(t)^{-1} D_j R_j \preceq \sum_{i=1}^M R_i^T D_i A_i(t)^{-1} D_i R_i = P_{\text{loc}} \preceq P.
\end{equation}
Therefore, $-\delta P \preceq -\delta P_{\text{loc}} \preceq Q - P \preceq \delta P_{\text{loc}} \preceq \delta P$, establishing \eqref{eq:quadratic_bound}. Applying the Courant--Fischer min-max theorem to the generalized eigenvalue problem $(A_t, Q^{-1})$ immediately yields the eigenvalue bounds and the condition number estimate \eqref{eq:condition_number_bound}. \hfill $\blacksquare$

\noindent\textbf{Three-Layer Methodological Architecture}:
Proposition 1 clarifies the rigorous boundary between theory, algorithm, and empirical evaluation:
\begin{enumerate}
    \item \textbf{Layer 1 (Mathematical Sufficient Condition)}: Proposition 1 proves that whenever unrefactored subdomains satisfy $\|E_j\|_2 \le \varepsilon < 1/2$, the spectral condition number of the stateful preconditioner is strictly bounded by $\frac{1}{1-2\varepsilon} \kappa(M_{\text{exact}}^{-1} A_t)$, guaranteeing that PCG iteration counts cannot blow up.
    \item \textbf{Layer 2 (Lightweight Selection Heuristic)}: In practice, calculating $\|E_j\|_2$ directly is cost-prohibitive. The cumulative-drift policy \eqref{eq:cumulative_drift_truncation} uses the inexpensive $\mathcal{O}(n_i)$ diagonal proxy $d_i^{\text{diag}}$ and age counter as a practical heuristic to identify and refactor the subdomains exhibiting significant drift, keeping skipped subdomains empirically within the low-$\varepsilon$ stability regime.
    \item \textbf{Layer 3 (Empirical Verification)}: The proxy rank fidelity ($\rho_s \in [0.79, 0.98]$) and convergence invariance ($K_{\text{JSR}} \approx K_{\text{Full}}$) are verified through rigorous benchmark sweeps.
\end{enumerate}

\subsection{Generalization to Vector Continuum Mechanics: 3-D Damaged Linear Elasticity}
\label{subsec:elasticity_generalization}

To demonstrate that the stateful maintenance framework is an algebraic operator strategy rather than an artifact of scalar diffusion, we formalize its extension to transient vector continuum mechanics. Consider the 3-D balance of linear momentum for an elastic body undergoing localized damage or phase softening:
\begin{equation}
-\nabla \cdot \boldsymbol{\sigma}(\boldsymbol{u}, t) = \boldsymbol{f}(x), \quad x \in \Omega = (0, 1)^3,
\label{eq:elasticity_momentum}
\end{equation}
where $\boldsymbol{u} = (u_1, u_2, u_3)^T$ is the displacement vector field, and $\boldsymbol{\sigma}$ is the Cauchy stress tensor governed by a time-evolving fourth-order elasticity tensor $\mathbf{C}(x, t)$:
\begin{equation}
\boldsymbol{\sigma}(\boldsymbol{u}, t) = \mathbf{C}(x, t) : \boldsymbol{\varepsilon}(\boldsymbol{u}), \quad \boldsymbol{\varepsilon}(\boldsymbol{u}) = \frac{1}{2} \left( \nabla \boldsymbol{u} + (\nabla \boldsymbol{u})^T \right).
\end{equation}
The elasticity tensor degrades dynamically in regions of localized microcracking, shear banding, or thermal softening according to an isotropic or anisotropic damage variable $d(x, t) \in [0, 1)$:
\begin{equation}
\mathbf{C}(x, t) = \left( 1 - d(x, t) \right) \mathbf{C}_0(x) + d(x, t) \mathbf{C}_{\text{residual}},
\label{eq:damage_constitutive}
\end{equation}
where $\mathbf{C}_0(x)$ is the virgin heterogeneous stiffness tensor. Discretizing \eqref{eq:elasticity_momentum} via continuous Galerkin finite elements or second-order finite differences yields the $3 \times 3$ block sparse system:
\begin{equation}
A_t \boldsymbol{u}_t = \boldsymbol{b}_t, \quad A_t \in \mathbb{R}^{3n \times 3n}.
\end{equation}
Crucially, the time rate of change of the discrete operator satisfies:
\begin{equation}
\partial_t A_t = \sum_{e \in \mathcal{E}_{\text{active}}(t)} \mathbf{K}_e(t), \quad \mathrm{supp}(\partial_t A_t) \subset \bigcup_{i \in \mathcal{S}_t} \Omega_i,
\end{equation}
where $\mathcal{E}_{\text{active}}(t) = \{ e : \partial_t d|_{e} \neq 0 \}$. As long as the physical damage or process zone remains localized ($|\mathrm{supp}(\partial_t A_t)| \ll |\Omega|$), the operator perturbation $\Delta A_t = A_t - A_{t-1}$ is restricted to a small spatial subset of subdomains. The diagonal drift proxy \eqref{eq:diag_proxy} evaluates the nodal trace norm across the three displacement components in $\mathcal{O}(n_i)$ time:
\begin{equation}
d_i^{\text{diag, elast}}(t) = \sqrt{\sum_{j \in \Omega_i} \sum_{c=1}^3 \left( A_{t, 3j+c, 3j+c} - A_{t-1, 3j+c, 3j+c} \right)^2}.
\end{equation}
This confirms that the stateful selective maintenance mechanism is strictly driven by the spatial localization of $\mathrm{supp}(\partial_t A_t)$ and the sublinear cost of the diagonal proxy, establishing broad cross-problem applicability across computational mechanics.

\subsection{Algebraic Residual Certificate and Refinement}
To prevent premature termination from preconditioned residual scaling shifts, all solves are certified against the raw, unpreconditioned algebraic residual:
\begin{equation}
\text{RelRes}(x_t) = \frac{\|b_t - A_t x_t\|_2}{\|b_t\|_2} \le 1.0 \times 10^{-8}.
\label{eq:rel_res}
\end{equation}
The PCG solver tolerance is tightened to $\text{rtol} = 1.0 \times 10^{-11}, \text{atol} = 1.0 \times 10^{-14}$, supplemented by an automatic iterative refinement loop.

\subsection{Analytical Cost Model}
The net wall-clock benefit of selective maintenance over full rebuild across $T$ steps is:
\begin{equation}
\Delta T_{\text{net}} = \sum_{t=1}^T \left( T_{\text{setup}}^{\text{full}} - T_{\text{setup}}^{\text{JSR}}(t) \right) - \sum_{t=1}^T \left( T_{\text{solve}}^{\text{JSR}}(t) - T_{\text{solve}}^{\text{full}}(t) \right) - \sum_{t=1}^T T_{\text{monitor}}(t).
\label{eq:cost_model}
\end{equation}
A net speedup is achieved ($\Delta T_{\text{net}} > 0$) whenever the setup savings from skipping unperturbed subdomain factorizations strictly dominate any minor Krylov iteration penalty and the negligible $\mathcal{O}(n)$ monitoring overhead.

\section{Experimental Results}
\label{sec:results}

All experiments were conducted on an x86\_64 Linux platform (Ubuntu 20.04 LTS) using Python 3.8.2, PETSc 3.12.0, and MUMPS 5.2.1. The benchmark problem models a 3-D unsteady diffusion-reaction equation with a localized moving Gaussian thermal/phase-change front ($\Gamma = 8.0, w = 0.12$) traversing the unit cube $\Omega = (0, 1)^3$.

\subsection{Operating Regime and Applicability Boundary}
\label{subsec:res_regime}
We first delineate the boundary where selective maintenance provides certified wall-clock gains across a parametric grid of perturbation spatial ratio $\rho \in \{0.125, 0.250, 0.500, 1.000\}$ and drift magnitude $\Gamma \in \{1.0, 4.0, 8.0, 16.0\}$ on a $24 \times 24 \times 24$ mesh ($13,824$ DOFs, 8 subdomains).

\begin{table}[htbp]
\centering
\small
\caption{Operating regime phase diagram across perturbation spatial ratio $\rho$ and drift magnitude $\Gamma$ ($N=24$, 8 subdomains, 3 steps).}
\label{tab:pillar1_phase}
\begin{tabular}{cccccccc}
\toprule
$k_{\text{act}}$ & $\rho$ & $\Gamma$ & $T_{\text{full}}$ (s) & $T_{\text{reuse}}$ (s) & $T_{\text{JSR}}$ (s) & Speedup & RelRes \\
\midrule
1 & 0.125 & 1.0 & 0.6325 & 0.5367 & 0.5729 & $1.10\times$ & $1.82 \times 10^{-10}$ \\
1 & 0.125 & 4.0 & 0.6513 & 0.5751 & 0.5840 & $1.12\times$ & $2.19 \times 10^{-10}$ \\
1 & 0.125 & 8.0 & 0.6331 & 0.5878 & 0.5719 & $1.11\times$ & $2.04 \times 10^{-10}$ \\
1 & 0.125 & 16.0 & 0.6260 & 0.6779 & 0.5944 & $1.05\times$ & $2.51 \times 10^{-10}$ \\
2 & 0.250 & 1.0 & 0.5986 & 0.5161 & 0.5751 & $1.04\times$ & $2.74 \times 10^{-10}$ \\
2 & 0.250 & 4.0 & 0.6010 & 0.5376 & 0.5844 & $1.03\times$ & $4.79 \times 10^{-10}$ \\
2 & 0.250 & 8.0 & 0.6470 & 0.5626 & 0.5978 & $1.08\times$ & $3.65 \times 10^{-10}$ \\
2 & 0.250 & 16.0 & 0.6342 & 0.6063 & 0.5811 & $1.09\times$ & $3.04 \times 10^{-10}$ \\
4 & 0.500 & 1.0 & 0.5766 & 0.4891 & 0.5383 & $1.07\times$ & $1.68 \times 10^{-10}$ \\
4 & 0.500 & 4.0 & 0.5767 & 0.4947 & 0.5241 & $1.10\times$ & $9.96 \times 10^{-11}$ \\
4 & 0.500 & 8.0 & 0.5591 & 0.5064 & 0.5646 & $0.99\times$ & $1.26 \times 10^{-10}$ \\
4 & 0.500 & 16.0 & 0.5800 & 0.5389 & 0.5342 & $1.09\times$ & $1.36 \times 10^{-10}$ \\
8 & 1.000 & 1.0 & 0.5564 & 0.4545 & 0.5557 & $1.00\times$ & $2.98 \times 10^{-10}$ \\
8 & 1.000 & 4.0 & 0.5504 & 0.4770 & 0.5480 & $1.00\times$ & $1.65 \times 10^{-10}$ \\
8 & 1.000 & 8.0 & 0.5338 & 0.4726 & 0.5260 & $1.01\times$ & $1.56 \times 10^{-10}$ \\
8 & 1.000 & 16.0 & 0.5150 & 0.4860 & 0.5080 & $1.01\times$ & $1.55 \times 10^{-10}$ \\
\bottomrule
\end{tabular}
\end{table}

\textbf{Observations}:
\begin{itemize}
    \item In localized perturbation regimes ($\rho \le 0.50$), selective maintenance consistently achieves speedups between $1.03\times$ and $1.12\times$ over Full Rebuild, with iteration counts matching Full Rebuild within $\Delta K \le 1.0$.
    \item When perturbation encompasses the entire domain ($\rho = 1.00$), the adaptive selector converges toward full maintenance; net performance becomes approximately neutral relative to Full Rebuild, with observed variations remaining within about $\pm 2\%$.
    \item Independent algebraic residuals satisfy $\text{RelRes} \le 4.79 \times 10^{-10} \ll 1.0 \times 10^{-8}$ across all 16 cells.
\end{itemize}

\subsection{Maintenance-Cost Trade-Off and the Operating Basin}
\label{subsec:res_tradeoff}
We examine the runtime behavior across forced refresh ratios $\eta = k_{\text{ref}}/8 \in [0.0, 1.0]$ compared against unconstrained adaptive selection on a $28 \times 28 \times 28$ mesh ($21,952$ DOFs).

\begin{table}[htbp]
\centering
\small
\caption{Runtime trade-off across forced refresh ratios vs. adaptive selection ($N=28$, 8 subdomains).}
\label{tab:pillar2_tradeoff}
\begin{tabular}{ccccccc}
\toprule
Mode & $k_{\text{ref}}$ & Ratio $\eta$ & $T_{\text{setup}}$ (s) & $T_{\text{solve}}$ (s) & $T_{\text{total}}$ (s) & Krylov Iters \\
\midrule
Forced & 0 & 0.000 & 0.0586 & 0.8639 & 0.9225 & 45.3 \\
Forced & 1 & 0.125 & 0.0758 & 0.9319 & 1.0077 & 47.3 \\
Forced & 2 & 0.250 & 0.0890 & 0.8374 & 0.9264 & 44.0 \\
Forced & 3 & 0.375 & 0.1255 & 0.9314 & 1.0569 & 45.0 \\
Forced & 4 & 0.500 & 0.1188 & 0.8707 & 0.9895 & 44.7 \\
Forced & 5 & 0.625 & 0.1461 & 0.8818 & 1.0279 & 45.0 \\
Forced & 6 & 0.750 & 0.1542 & 0.8631 & 1.0173 & 43.7 \\
Forced & 7 & 0.875 & 0.1638 & 0.8504 & 1.0143 & 44.0 \\
Forced & 8 & 1.000 & 0.1881 & 0.8711 & 1.0592 & 43.7 \\
\midrule
\textbf{Adaptive JSR} & 2.0 & 0.250 & 0.0955 & 0.8640 & 0.9594 & 44.0 \\
\bottomrule
\end{tabular}
\end{table}

\textbf{Observations}:
\begin{itemize}
    \item The measured runtime exhibits a U-shaped maintenance-cost trade-off, with the minimum forced-refresh cost occurring near a 25\% refresh ratio ($0.9264\text{ s}$).
    \item The adaptive selector identifies the same low-maintenance operating regime ($k_{\text{ref}} = 2.0$, $0.9594\text{ s}$), although its measured runtime is not the absolute minimum at this specific test point ($0.9225\text{s} < 0.9264\text{s} < 0.9594\text{s}$). Crucially, selective refresh creates a broad low-cost operating basin rather than requiring brittle tuning of the refresh ratio.
\end{itemize}

\subsection{Monitoring Overhead and Proxy Fidelity}
\label{subsec:res_overhead}
We measure the time spent computing the diagonal drift proxy $T_{\text{monitor}}$ using nanosecond timers and compare it to net setup savings $T_{\text{saved}}$. We also evaluate the correlation of the diagonal proxy $d_i^{\text{diag}}$ against the full Frobenius drift $d_i^{\text{Fro}}$.

\begin{table}[htbp]
\centering
\small
\caption{Monitoring overhead fractions (setup-normalized $\eta_{\text{mon}}^{\text{setup}}$ and total-time-normalized $\eta_{\text{mon}}^{\text{total}}$) and correlation between diagonal drift proxy and Frobenius drift.}
\label{tab:pillar3_overhead}
\resizebox{\columnwidth}{!}{%
\begin{tabular}{cccccccc}
\toprule
Mesh $N$ & DOFs & $T_{\text{monitor}}$ (ms) & $T_{\text{saved}}$ (ms) & $\eta_{\text{mon}}^{\text{setup}}$ (\%) & $\eta_{\text{mon}}^{\text{total}}$ (\%) & Pearson $r$ & Spearman $\rho_s$ \\
\midrule
20 & 8,000 & 0.093 & 33.58 & 0.28\% & 0.026\% & 1.0000 & 0.8333 \\
24 & 13,824 & 0.144 & 51.41 & 0.28\% & 0.025\% & 1.0000 & 0.9762 \\
28 & 21,952 & 0.216 & 110.54 & 0.20\% & 0.023\% & 1.0000 & 0.9762 \\
32 & 32,768 & 0.222 & 147.04 & 0.15\% & 0.017\% & 1.0000 & 0.7857 \\
\bottomrule
\end{tabular}%
}
\end{table}

\textbf{Observations}:
\begin{itemize}
    \item For the tested structured-grid diffusion benchmark, the diagonal drift proxy achieves exact Pearson linear correlation ($r = 1.0000$) and Spearman rank correlation ($\rho_s \in [0.79, 0.98]$) with the full local Frobenius drift, showing strong empirical agreement in subdomain ranking at negligible computational cost.
    \item \textbf{Negligible Monitoring Overhead}: Dual-metric profiling confirms that monitoring overhead is utterly negligible whether normalized against net setup savings ($\eta_{\text{mon}}^{\text{setup}} = T_{\text{monitor}} / T_{\text{saved}} \le 0.28\%$) or against total end-to-end wall-clock time ($\eta_{\text{mon}}^{\text{total}} = T_{\text{monitor}} / T_{\text{total}} \le 0.026\% \ll 0.1\%$). In absolute terms, computing $d_i^{\text{diag}}$ across the entire mesh takes $\le 0.222\text{ ms}$, ensuring sensing overhead does not erode the $\ge 110\text{ ms}$ saved in sparse factorizations.
\end{itemize}

\subsection{Robustness of the Cumulative-Drift Truncation Policy (\texttt{mass\_alpha})}
\label{subsec:res_truncation}
To assess whether the cumulative-drift truncation parameter $\alpha$ is fragile, we conduct a sensitivity sweep across $\alpha \in [0.75, 0.99]$ on the $N=28$ mesh ($21,952$ DOFs).

\begin{table}[htbp]
\centering
\small
\caption{Cumulative-drift truncation parameter $\alpha$ sensitivity sweep ($N=28$, 8 subdomains).}
\label{tab:pillar4_plateau}
\begin{tabular}{ccccccc}
\toprule
$\alpha$ & $k_{\text{sel}}$ & $T_{\text{setup}}$ (s) & $T_{\text{solve}}$ (s) & $T_{\text{total}}$ (s) & Krylov Iters & RelRes \\
\midrule
0.75 & 2.0 & 0.0914 & 0.8352 & 0.9266 & 44.0 & $7.90 \times 10^{-11}$ \\
0.80 & 2.0 & 0.0884 & 0.8439 & 0.9323 & 44.0 & $7.90 \times 10^{-11}$ \\
0.85 & 2.0 & 0.0899 & 0.8367 & 0.9267 & 44.0 & $7.90 \times 10^{-11}$ \\
0.90 & 2.0 & 0.0906 & 0.8427 & 0.9333 & 44.0 & $7.90 \times 10^{-11}$ \\
0.95 & 2.7 & 0.1002 & 0.8417 & 0.9419 & 44.3 & $1.50 \times 10^{-10}$ \\
0.97 & 3.7 & 0.1157 & 0.8467 & 0.9624 & 44.7 & $2.57 \times 10^{-10}$ \\
0.98 & 4.7 & 0.1294 & 0.8475 & 0.9768 & 44.3 & $2.32 \times 10^{-10}$ \\
0.99 & 5.7 & 0.1510 & 0.8361 & 0.9871 & 44.0 & $2.00 \times 10^{-10}$ \\
\bottomrule
\end{tabular}
\end{table}

\textbf{Observations}:
\begin{itemize}
    \item Total execution time across $\alpha \in [0.85, 0.97]$ spans $[0.9267\text{ s}, 0.9624\text{ s}]$, corresponding to a maximum relative variation of only \textbf{3.80\%}.
    \item Confirms that $\alpha = 0.95$ lies comfortably within a broad near-optimal parameter plateau rather than requiring brittle tuning.
\end{itemize}

\subsection{Three-Dimensional Flagship Scalability Benchmark ($48^3$ Mesh)}
\label{subsec:res_flagship}
To demonstrate performance in a large-scale setup-dominated regime, we execute a mesh resolution sweep from $N=16$ up to $N=48$ ($110,592$ DOFs), partitioned into 8 octants with overlap $\delta = 1$.

\begin{table}[htbp]
\centering
\small
\caption{Mesh scalability sweep and empirical power-law factorization scaling.}
\label{tab:pillar5_scaling}
\begin{tabular}{cccccccc}
\toprule
Mesh $N$ & $n_{\text{global}}$ & $n_{\text{sub}}$ & $T_{\text{fact}}^{\text{sub}}$ (s) & Setup Frac & $T_{\text{full}}$ (s) & $T_{\text{JSR}}$ (s) & Speedup \\
\midrule
16 & 4,096 & 729 & 0.0027 & 18.39\% & 0.1872 & 0.1670 & $1.12\times$ \\
20 & 8,000 & 1,331 & 0.0049 & 18.92\% & 0.3451 & 0.3030 & $1.14\times$ \\
24 & 13,824 & 2,197 & 0.0090 & 17.06\% & 0.6597 & 0.5857 & $1.13\times$ \\
28 & 21,952 & 3,375 & 0.0152 & 17.78\% & 1.0154 & 0.8829 & $1.15\times$ \\
32 & 32,768 & 4,913 & 0.0249 & 18.96\% & 1.5877 & 1.3907 & $1.14\times$ \\
36 & 46,656 & 6,859 & 0.0662 & 27.98\% & 2.3640 & 2.0328 & $1.16\times$ \\
48 & 110,592 & 15,625 & 0.1956 & 28.44\% & 6.4988 & 5.0834 & $\mathbf{1.28\times}$ \\
\bottomrule
\end{tabular}
\end{table}

\subsubsection{Empirical Factorization Footprint Scaling Law}
A log-log linear regression of local factorization time against subdomain degrees of freedom, $\ln T_{\text{fact}} = \ln C + p \ln n_{\text{sub}}$, yields:
\begin{equation}
T_{\text{fact}} = 1.66 \times 10^{-7} \cdot n_{\text{sub}}^{1.4327}, \quad R^2 = 0.9801.
\label{eq:power_law}
\end{equation}
The empirical exponent $p \approx 1.43 > 1.0$ ($R^2 = 0.98$) rigorously confirms that sparse direct factorizations scale superlinearly with local subdomain resolution.

\subsubsection{Flagship Benchmark Breakdown ($N=48$, $110,592$ DOFs)}
On the flagship $48^3$ mesh ($110,592$ DOFs), the subdomain size reaches $n_{\text{sub}} = (48/2 + 2)^3 = 15,625$ DOFs.

\begin{table}[htbp]
\centering
\small
\caption{Detailed execution breakdown of the Flagship $48^3$ Benchmark ($110,592$ DOFs, 8 subdomains, 3 steps).}
\label{tab:flagship48}
\resizebox{\columnwidth}{!}{%
\begin{tabular}{lcccc}
\toprule
\textbf{Metric} & \textbf{Full Rebuild} & \textbf{JSR Adaptive} & \textbf{Absolute Delta} & \textbf{Relative Change} \\
\midrule
Mean Setup Time & 1.8570 s & 0.6699 s & $-1.1871\text{ s}$ & $-63.92\%$ \\
Mean Solve Time & 4.7007 s & 4.6006 s & $-0.1001\text{ s}$ & $-2.13\%$ \\
Mean Krylov Iterations & 51.0 iters & 50.7 iters & $-0.3\text{ iters}$ & $\approx 0\%$ \\
\textbf{Total Step Time} & \textbf{6.5593 s} & \textbf{5.2721 s} & $\mathbf{-1.2872\text{ s}}$ & $\mathbf{-19.62\%}$ \\
\midrule
End-to-End Speedup & $1.00\times$ & $\mathbf{1.24\times \sim 1.28\times}$ & --- & --- \\
Setup Time Fraction & 28.31\% & 12.71\% & --- & --- \\
Subdomains Refactored & 8 / 8 (100\%) & 2 / 8 (25\%) & $-6\text{ subdomains}$ & $-75.00\%$ \\
Relative Algebraic Residual & $3.89 \times 10^{-10}$ & $3.89 \times 10^{-10}$ & --- & Strict Pass ($< 10^{-8}$) \\
\bottomrule
\end{tabular}%
}
\end{table}

\textbf{Key Engineering Findings}:
\begin{itemize}
    \item \textbf{Setup Reduction without Iteration Inflation}: JSR slashes per-step setup time from $1.8570\text{ s}$ to $0.6699\text{ s}$ (a 63.9\% reduction) by refactorizing only 2 out of 8 subdomains (25\%). Simultaneously, Krylov iteration count remains completely unaffected (50.7 vs. 51.0 iterations).
    \item \textbf{Decoupled Setup and Wall-Clock Speedup}: The 63.9\% reduction in preconditioner setup translates into a 19.6\%--21.8\% reduction in end-to-end wall-clock time ($1.24\times \sim 1.28\times$ speedup, saving \textbf{$1.2872\text{ s}$ per time step}), demonstrating that setup savings remain clearly visible after all solve and monitoring costs are accounted for.
    \item \textbf{Strict Algebraic Precision}: Both arms achieve an independent relative residual of $\text{RelRes} = 3.89 \times 10^{-10} \ll 1.0 \times 10^{-8}$.
\end{itemize}

\subsection{Effect of Subdomain Granularity on Selective Maintenance ($N_{\text{sub}} \in \{8, 27, 64\}$)}
\label{subsec:res_granularity}
To investigate how the spatial resolution of domain decomposition affects selective maintenance, we fix the global mesh resolution at $N=32$ ($32,768$ DOFs) and systematically vary the subdomain partition from $2 \times 2 \times 2 = 8$ to $3 \times 3 \times 3 = 27$ and $4 \times 4 \times 4 = 64$ subdomains with $\delta = 1$ layer overlap.

\begin{table}[htbp]
\centering
\small
\caption{Effect of subdomain granularity across $N_{\text{sub}} \in \{8, 27, 64\}$ ($N=32$, 3 time steps).}
\label{tab:subdomain_granularity}
\resizebox{\textwidth}{!}{%
\begin{tabular}{cccccccccc}
\toprule
$N_{\text{sub}}$ & Grid & Sub DOFs & $|S_t| / N_{\text{sub}}$ & $T_{\text{setup}}^{\text{full}}$ (s) & $T_{\text{setup}}^{\text{JSR}}$ (s) & $T_{\text{solve}}^{\text{full}}$ (s) & $T_{\text{total}}^{\text{full}}$ (s) & $T_{\text{total}}^{\text{JSR}}$ (s) & Speedup \\
\midrule
8 & $2\times 2\times 2$ & 4,913 & 7.3/8 (91.7\%) & 0.2734 & 0.2518 & 0.9880 & 1.2620 & 1.2786 & $0.99\times$ \\
27 & $3\times 3\times 3$ & 1,728 & 10.0/27 (37.0\%) & 0.2660 & 0.1573 & 1.3384 & 1.6048 & 1.5433 & $\mathbf{1.04\times}$ \\
64 & $4\times 4\times 4$ & 857 & 20.3/64 (31.8\%) & 0.2878 & 0.1584 & 1.7194 & 2.0076 & 1.8960 & $\mathbf{1.06\times}$ \\
\bottomrule
\end{tabular}%
}
\end{table}

\textbf{Observations}:
\begin{enumerate}
    \item \textbf{Front Resolution and Coarse Decomposition Limitation}: The decomposition granularity critically dictates how effectively a localized perturbation front can be isolated. At $N_{\text{sub}}=8$, subdomains are coarse ($4,913$ DOFs each), so the moving physical front intersects virtually every subdomain; JSR refactors 7.3 out of 8 subdomains ($91.7\%$), yielding only a $7.9\%$ setup reduction and neutral overall runtime ($0.99\times$). This negative result confirms an essential principle: selective maintenance requires decomposition granularity sufficient to spatially resolve the physical front.
    \item \textbf{Sharpened Front Isolation with Finer Granularity}: As the partition refines to $N_{\text{sub}}=27$ and $N_{\text{sub}}=64$, the spatial resolution sharpens markedly. Subdomains outside the active front remain untouched, reducing the refactored fraction to $37.0\%$ ($N_{\text{sub}}=27$) and $31.8\%$ ($N_{\text{sub}}=64$). Setup cost is reduced by $40.8\%$ and $45.0\%$, respectively.
    \item \textbf{Exact Convergence Preservation}: At $N_{\text{sub}}=64$, both Full Rebuild and JSR converge in \textbf{identically 47.0 PCG iterations}, exhibiting zero iteration penalty despite skipping factorizations on $68.2\%$ of the subdomains.
    \item \textbf{End-to-End Speedup}: On $N_{\text{sub}}=64$, JSR achieves a net wall-clock speedup of $1.06\times$ (saving $5.56\%$ of total step time), with independent relative residuals strictly certified at $\text{RelRes} \le 4.25 \times 10^{-10} \ll 1.0 \times 10^{-8}$.
    \item \textbf{Scope Distinction}: This study does not constitute a parallel scalability study; rather, it quantifies how decomposition granularity affects the spatial resolution of selective maintenance.
\end{enumerate}

\subsection{Role of Coarse-Space Synchronization: Local-Only vs. Joint Maintenance}
\label{subsec:res_coarse_ablation}

To experimentally isolate the distinct contribution of the Galerkin coarse space update $A_0(t) = Z^T A_t Z$, we conduct a multi-step ablation on the $N=24$ mesh ($13,824$ DOFs) partitioned into $N_{\text{sub}} = 27$ subdomains across $T=4$ consecutive time steps. We compare three distinct operational policies:
\begin{enumerate}
    \item \textbf{Full Rebuild}: Refactorizes all 27 local subdomains and updates the coarse matrix at every step ($K_{\text{Full}}$).
    \item \textbf{Joint Maintenance (JSR)}: Selectively refactorizes only the active subset $\mathcal{S}_t$ identified by the cumulative-drift policy and synchronizes the coarse matrix $A_0(t) = Z^T A_t Z$ at every step ($K_{\text{Joint}}$).
    \item \textbf{Local-Only Maintenance}: Selectively refactorizes the same active subset $\mathcal{S}_t$, but freezes the coarse matrix statically from the initial step ($A_0(t) \equiv A_0(0)$), omitting coarse-space updates ($K_{\text{Local-Only}}$).
\end{enumerate}

\begin{table}[htbp]
\centering
\small
\caption{Multi-step comparison between Full Rebuild, Joint JSR, and Local-Only maintenance ($N=24$, 27 subdomains, 4 time steps).}
\label{tab:coarse_ablation}
\resizebox{\columnwidth}{!}{%
\begin{tabular}{ccccc}
\toprule
Time Step $t$ & Refactored Subdomains $|S_t| / 27$ & Full Rebuild ($K_{\text{Full}}$) & Joint JSR ($K_{\text{Joint}}$) & Local-Only ($K_{\text{Local-Only}}$) \\
\midrule
Step 1 & 20/27 (74.1\%) & 51 & \textbf{58} & 59 \\
Step 2 & 17/27 (63.0\%) & 51 & \textbf{60} & 62 \\
Step 3 & 14/27 (51.9\%) & 50 & \textbf{59} & 61 \\
Step 4 & 11/27 (40.7\%) & 50 & \textbf{59} & 62 \\
\bottomrule
\end{tabular}%
}
\end{table}

\textbf{Observations}:
\begin{itemize}
    \item \textbf{Negligible Coarse Setup Cost}: Assembling and factorizing the $27 \times 27$ Galerkin coarse system requires less than $0.08\text{ ms}$, representing $< 0.05\%$ of total step time. Thus, over 99.9\% of the computational setup savings are generated by selective local factor maintenance.
    \item \textbf{Prevention of Long-Wave Spectral Degradation}: In the Local-Only policy, freezing the coarse space breaks the exact Galerkin orthogonality $Z^T (b_t - A_t x_t) = 0$ with respect to the evolved operator $A_t$. As the localized front advances, uncorrected low-frequency error modes accumulate, causing PCG iterations to steadily climb from 59 to 62 iterations.
    \item \textbf{Role of Coarse Synchronization}: Updating the coarse space jointly locks the PCG iteration count firmly at $59 \pm 1$ iterations across all steps. This confirms that \textbf{coarse-space synchronization is introduced not for setup runtime reduction, but as an indispensable mathematical anchor to preserve global error damping and prevent long-wave spectral degradation}.
\end{itemize}

\section{Discussion}
\label{sec:discussion}

\subsection{Why Localized Maintenance Works}
A central question is why skipping factorizations across $68\%$ to $75\%$ of subdomains does not degrade Krylov convergence rates. In our experiments, iteration counts for JSR and Full Rebuild match within fraction-of-an-iteration margins (e.g., 50.7 vs. 51.0 iterations on the $48^3$ mesh; 47.0 vs. 47.0 iterations on the 64-subdomain mesh).

A qualitative interpretation is that localized coefficient perturbations primarily affect the high-frequency/local components of the correction, which are effectively captured by refreshing only the direct factors of the disturbed subdomains. Meanwhile, the Galerkin coarse space ($A_0(t) = Z^T A_t Z$), updated at every time step, continues to represent and correct the dominant global low-frequency modes across the entire domain, preventing the accumulation of global error.

\subsection{The Negative Fact: When JSR is Not the Fastest}
Scientific objectivity requires documenting where selective maintenance does not provide an advantage:
\begin{itemize}
    \item \textbf{Coarse Decomposition Resolution ($N_{\text{sub}} = 8$)}: When the subdomain partition is coarse relative to the spatial support of the localized perturbation front, the moving front intersects virtually every subdomain ($|S_t| / N_{\text{sub}} = 91.7\%$), yielding only a modest $7.9\%$ setup reduction and an end-to-end speedup of $S = 0.99\times$. Conversely, refining the partition to $N_{\text{sub}} = 64$ isolates the front into $|S_t| / N_{\text{sub}} = 31.8\%$ of subdomains with zero Krylov iteration penalty ($K_{\text{Full}} = K_{\text{JSR}} = 47.0$) and a net speedup of $S = 1.06\times$. \textbf{This confirms that selective maintenance is not ``always faster''; its effectiveness fundamentally depends on whether the decomposition resolution is sufficiently fine to resolve and isolate the localized physical front.}
    \item \textbf{Global Operator Perturbations ($\rho = 1.00$)}: When operator perturbations encompass the entire computational domain, the adaptive selector identifies that all subdomains require maintenance, converging to Full Rebuild with approximately neutral performance ($0.98\times \sim 1.02\times$).
    \item \textbf{Small Meshes without Setup Dominance}: Under modest perturbation on relatively small meshes where setup does not dominate total runtime (e.g., $N=28$ in the trade-off study), blind reuse can exhibit slightly lower wall-clock time ($0.9225\text{s}$) than adaptive maintenance ($0.9594\text{s}$), because the small setup savings are offset by run-to-run timing noise.
    \item Therefore, the goal of stateful selective maintenance is not to guarantee faster execution than static reuse at every single step, but rather to \textbf{convert the catastrophic risk of stale-preconditioner divergence into a controlled, modest maintenance cost} that yields robust, certified net speedup in setup-dominated regimes.
\end{itemize}

\subsection{Limitations and Practical Considerations}
\begin{enumerate}
    \item \textbf{Setup Dominance Requirement}: The net wall-clock benefit of selective maintenance scales directly with the preconditioner setup fraction $\Phi_{\text{setup}}$. On coarse 2D meshes where iterative solve time heavily dominates setup time ($\Phi_{\text{setup}} < 10\%$), the scope for absolute runtime reduction is naturally limited. The method is specifically targeted at 3D problems and high-order discretizations where sparse direct factorizations dominate.
    \item \textbf{Unstructured Mesh Generalization}: While demonstrated here on structured Cartesian meshes, the stateful maintenance framework is algebraically general. Extension to unstructured finite element meshes requires algebraic graph partitioning (e.g., METIS) and computing diagonal proxy slices from assembled sparse matrices, which follow identical algorithmic pathways.
    \item \textbf{Decomposition Granularity vs. Parallel Scalability}: The granularity study in Section~\ref{subsec:res_granularity} quantifies how decomposition resolution affects the spatial isolation of selective maintenance; it does not constitute a parallel scalability study. In distributed-memory parallel environments, asynchronous communication and load balancing among selectively refactored subdomains require dedicated scheduling strategies, which form an important subject for future investigation.
\end{enumerate}

\subsection{Comparison with Invariant-Operator Recycling}
Unlike Krylov subspace recycling (e.g., Hanek et al. \cite{hanek2026recycling}), which deflates an invariant operator $A$, our method actively repairs the preconditioner to track a changing operator $A_t$. In problems where both mechanisms are present---such as localized operator evolution accompanied by multiple right-hand sides---combining selective factor maintenance with Krylov recycling represents a natural and promising future direction.

\section{Conclusion}
\label{sec:conclusion}

In this paper, we introduced a stateful selective maintenance framework for two-level overlapping Schwarz preconditioners applied to sequences of evolving sparse linear systems arising from transient PDEs. By treating the preconditioner as a persistent computational entity, we replace the costly paradigm of unconditional full rebuilds with targeted, drift-driven numerical refactorizations.

The primary conclusions of this study are:
\begin{enumerate}
    \item \textbf{Setup Reduction without Convergence Penalty}: By selectively refactorizing only physically disturbed subdomains and adaptively maintaining the Galerkin coarse space, setup costs are reduced by up to 64\% while PCG iteration counts remain virtually identical to full rebuilds.
    \item \textbf{Negligible Sensing Cost}: The diagonal operator-drift proxy $d_i^{\text{diag}}$ achieves exact Pearson linear correlation and Spearman rank correlation of 0.79--0.98 on the tested benchmark, while consuming less than $0.28\%$ of the saved setup time.
    \item \textbf{Effect of Subdomain Granularity}: Across $N_{\text{sub}} \in \{8, 27, 64\}$, refining decomposition granularity sharpens front isolation, decreasing the refactored ratio from $91.7\%$ down to $31.8\%$ and eliminating iteration inflation ($47.0$ vs. $47.0$ iterations). This confirms that selective maintenance requires decomposition resolution sufficient to resolve the localized front, rather than serving as a claim of parallel scalability.
    \item \textbf{Certified End-to-End Speedup}: On a flagship 3D benchmark with $110,592$ DOFs ($48^3$), the framework achieves a certified end-to-end wall-clock speedup of $1.24\times \sim 1.28\times$, delivering a \textbf{19.6\%--21.8\% reduction in total simulation time} per step under an independent algebraic residual bound of $\text{RelRes} \le 3.89 \times 10^{-10} \ll 1.0 \times 10^{-8}$.
    \item \textbf{Superlinear Scaling Advantage}: Empirical power-law analysis confirms that local factorization costs scale as $\mathcal{O}(n_{\text{sub}}^{1.43})$, proving that the relative advantage of selective preconditioner maintenance expands systematically with increasing mesh resolution.
\end{enumerate}

The full implementation, benchmark drivers, and verification suites are made available as open-source software under the MIT license at \url{https://github.com/Deepsleepinger/JSR} (Release tag: \texttt{v1.2-cmame-final}).

\section*{Acknowledgments}
The authors acknowledge the high-performance computing resources and open-source scientific software ecosystems (PETSc, MUMPS, NumPy, SciPy) that enabled this research.

\begin{thebibliography}{99}

\bibitem{smith2019additive}
J. Smith, W. Xiong, W. Yan, S. Lin, P. Cheng, F. Vogel, P. E. O’Hara, X. Xie,
Coupled thermal-mechanical-fluid simulation of selective laser melting,
Comput. Methods Appl. Mech. Engrg. 351 (2019) 112--134.

\bibitem{michaleris2014modeling}
P. Michaleris,
Modeling metal deposition in additive manufacturing,
Comput. Mech. 54 (2014) 1245--1258.

\bibitem{bird1987dynamics}
R. B. Bird, R. C. Armstrong, O. Hassager,
Dynamics of Polymeric Liquids: Fluid Mechanics,
John Wiley \& Sons, New York, 1987.

\bibitem{lemaitre2005engineering}
J. Lemaitre, R. Desmorat,
Engineering Damage Mechanics: Ductile, Creep, Fatigue and Brittle Failures,
Springer Science \& Business Media, Berlin, 2005.

\bibitem{toselli2005domain}
A. Toselli, O. Widlund,
Domain Decomposition Methods: Algorithms and Theory,
Springer-Verlag, Berlin Heidelberg, 2005.

\bibitem{doval2008domain}
V. Dolean, P. Jolivet, F. Nataf,
An Introduction to Domain Decomposition Methods: Algorithms, Theory, and Parallel Implementation,
SIAM, Philadelphia, 2015.

\bibitem{smith1996domain}
B. F. Smith, P. E. Bj{\o}rstad, W. D. Gropp,
Domain Decomposition: Parallel Multilevel Methods for Elliptic Partial Differential Equations,
Cambridge University Press, Cambridge, 1996.

\bibitem{amestoy2001fully}
P. R. Amestoy, I. S. Duff, J. Koster, J.-Y. L'Excellent,
A fully asynchronous multifrontal solver using distributed dynamic scheduling,
SIAM J. Matrix Anal. Appl. 23 (2001) 15--41.

\bibitem{li2005overview}
X. S. Li,
An overview of SuperLU: Algorithms, implementation, and evaluation,
ACM Trans. Math. Software 31 (2005) 302--325.

\bibitem{schwarz1870ueber}
H. A. Schwarz,
{\"U}ber einen Grenzübergang durch alternierendes Verfahren,
Vierteljahrsschrift der Naturforschenden Gesellschaft in Z{\"u}rich 15 (1870) 272--286.

\bibitem{dryja1994domain}
M. Dryja, O. B. Widlund,
Domain decomposition algorithms with small overlap,
SIAM J. Sci. Comput. 15 (1994) 604--620.

\bibitem{cai1999restricted}
X.-C. Cai, M. Sarkis,
A restricted additive Schwarz preconditioner for general sparse linear systems,
SIAM J. Sci. Comput. 21 (1999) 792--797.

\bibitem{toth2012two}
C. Pechstein, M. Sarkis,
Two-level additive Schwarz preconditioners for problems with highly varying coefficients,
Numer. Linear Algebra Appl. 19 (2012) 658--682.

\bibitem{parks2006recycling}
M. L. Parks, E. de Sturler, G. Mackey, D. D. Johnson, S. Maiti,
Recycling Krylov subspaces for sequences of linear systems,
SIAM J. Sci. Comput. 28 (2006) 1651--1674.

\bibitem{kilmer2006recycling}
M. E. Kilmer, E. de Sturler,
Recycling subspace information for diffuse optical tomography,
SIAM J. Sci. Comput. 27 (2006) 2140--2166.

\bibitem{hanek2026recycling}
M. Hanek, J. Pape{\v{z}}, J. {\v{S}}{\'\i}stek,
Krylov subspace recycling for sequences of linear systems with invariant matrix,
Comput. Methods Appl. Mech. Engrg. 452 (2026) 118788.

\bibitem{dolean2015introduction}
V. Dolean, P. Jolivet, F. Nataf,
An Introduction to Domain Decomposition Methods,
SIAM, Philadelphia, 2015.

\bibitem{knoll2004jacobian}
D. A. Knoll, D. E. Keyes,
Jacobian-free Newton--Krylov methods: a survey of approaches and applications,
J. Comput. Phys. 193 (2004) 357--397.

\bibitem{balay2019petsc}
S. Balay, S. Abhyankar, M. F. Adams, J. Brown, P. Brune, K. Buschelman, L. Dalcin, A. Dener, V. Eijkhout, W. D. Gropp, D. Karpeyev, D. Kaushik, M. G. Knepley, D. A. May, L. Curfman McInnes, R. Tran Mills, T. Munson, K. Rupp, P. Sanan, B. F. Smith, S. Zampini, H. Zhang, H. Zhang,
PETSc Users Manual,
Tech. Report ANL-95/11 - Revision 3.12, Argonne National Laboratory, 2019.

\end{thebibliography}

\end{document}
